\documentclass[a4paper,,onecolumn,10pt]{article}
\usepackage[spanish]{babel}
\usepackage[dvipsnames]{xcolor}
\usepackage[latin1]{inputenc}
\usepackage{listings}

\lstdefinestyle{estilocodigo}{
backgroundcolor=\color{gray!12},
keywordstyle=\color{magenta},
stringstyle=\color{blue!50!red},
showstringspaces=false,
numbers=left,
frame=single,
breaklines=true
}

\begin{document}
\pagecolor{gray} %color fondo gris
\color{blue} %letras en azul

\begin{lstlisting}[language=C,basicstyle=\footnotesize, style=estilocodigo]
#include<pthread.h>
#include<stdio.h>
#include<stdlib.h>

int main (int argc , char * argv[]){
   int valor=20;
   int * puntero=NULL;
   puntero = &valor;
   *puntero = 40;
   printf("Direccion de memoria de valor: %p\n", &valor);
   printf("Direccion de memoria de puntero: %p\n", &puntero);
   printf("Puntero numero: %d\n",*puntero);
   printf("Valor numero: %d\n",valor);
   
}
\end{lstlisting}

\textbf{Solucion}

Direccion de memoria de valor:0x7ffead1bb5f4

Direccion de memoria de valor:0x7ffead1bb5f8

Puntero numero:40

Valor numero:40
\vspace{7.5mm}

Valor = 20 

*Puntero = NULL

*Puntero $\rightarrow$ NULL

Puntero $\rightarrow$ 0x7ffead1bb5f4 (20)

*Puntero = 40

Puntero es 0x7ffead1bb5f8(40) 

Puntero $\rightarrow$ 0x7ffead1bb5f4 (40)

Valor = 40 

\vspace{10mm}

{\color{RubineRed} \rule{\linewidth}{0.5mm} }


\begin{lstlisting}[language=C,basicstyle=\footnotesize, style=estilocodigo]
#include<pthread.h>
#include<stdio.h>
#include<stdlib.h>

int main (int argc , char * argv[]){
   int valor=20;
   int * puntero=NULL;
   puntero = &valor;
   puntero = 40;
   printf("Direccion de memoria de valor: %p\n", &valor);
   printf("Direccion de memoria de puntero: %p\n", &puntero);
   printf("Puntero numero: %d\n",puntero);
   printf("Valor numero: %d\n",valor);
   
}
\end{lstlisting}

\textbf{Solucion}

Direccion de memoria de valor:0x7ffead1bb5f4

Direccion de memoria de valor:0x7ffead1bb5f8

Puntero numero:40

Valor numero:20
\vspace{7.5mm}

Valor = 20 

*Puntero = NULL

*Puntero $\rightarrow$ NULL

Puntero $\rightarrow$ 0x7ffead1bb5f4 (20)

Puntero = 40

Puntero es 0x7ffead1bb5f8(40) 

Puntero $\rightarrow$ 0x7ffead1bb5f4 (20)

Valor = 20

\vspace{10mm}

{\color{RubineRed} \rule{\linewidth}{0.5mm} }

\begin{lstlisting}[language=C,basicstyle=\footnotesize, style=estilocodigo]
#include<pthread.h>
#include<stdio.h>
#include<stdlib.h>

void valorsuma(int * mipuntero){
    int suma= * mipuntero+20;
    printf("Valor de suma: %d\n",suma);
}

int main (int argc , char * argv[]){
   int * puntero=NULL;
   puntero=40;
   printf("Puntero numero: %d\n",puntero);
   valorsuma(&puntero);
   printf("Puntero numero: %d\n",puntero);
   
}
\end{lstlisting}

\textbf{Solucion}

Puntero numero:40

Valor de suma: 60 

Puntero numero:40

\vspace{7.5mm}

*Puntero = NULL

*Puntero $\rightarrow$ NULL

Puntero = 40

int *mipuntero = $\&$puntero

*mipuntero $\rightarrow$ direccion puntero(40)

mipuntero=40

\vspace{10mm}

{\color{RubineRed} \rule{\linewidth}{0.5mm} }

\begin{lstlisting}[language=C,basicstyle=\footnotesize, style=estilocodigo]
#include<pthread.h>
#include<stdio.h>
#include<stdlib.h>

void valorsuma(int * mipuntero){
    * mipuntero= * mipuntero+20;
    printf("Valor de suma: %d\n",* mipuntero);
}

int main (int argc , char * argv[]){
   int * puntero=NULL;
   puntero=40;
   printf("Puntero numero: %d\n",puntero);
   valorsuma(&puntero);
   printf("Puntero numero: %d\n",puntero);
   
}
\end{lstlisting}

\textbf{Solucion}

Puntero numero:40

Valor de suma: 60 

Puntero numero:60

\vspace{7.5mm}

*Puntero = NULL

*Puntero $\rightarrow$ NULL

Puntero = 40

int *mipuntero = $\&$puntero

*mipuntero $\rightarrow$ direccion puntero(40)

mipuntero=40

*mipuntero= *mipuntero $\rightarrow$ direccion puntero(40) + 20

*mipuntero $\rightarrow$ direccion puntero(60)

Puntero=60

\vspace{10mm}

{\color{RubineRed} \rule{\linewidth}{0.5mm} }

\begin{lstlisting}[language=C,basicstyle=\footnotesize, style=estilocodigo]
#include<pthread.h>
#include<stdio.h>
#include<stdlib.h>

void valorsuma(int * mipuntero){
    * mipuntero= * mipuntero+20;
    printf("Valor de suma: %d\n",* mipuntero);
}

int main (int argc , char * argv[]){
   int numero=40;
   printf("Numero: %d\n",numero);
   valorsuma(&numero);
   printf("Numero: %d\n",numero);
   
}
\end{lstlisting}

\textbf{Solucion}

Numero:40

Valor de suma: 60 

Numero:60

\vspace{7.5mm}

Numero = 40

int *mipuntero = $\&$numero

*mipuntero $\rightarrow$ direccion numero(40)

mipuntero=40

*mipuntero= *mipuntero $\rightarrow$ direccion puntero(40) + 20

*mipuntero $\rightarrow$ direccion numero(60)

Numero=60

\vspace{10mm}

\end{document}